۳۱۴
۳۱۵
«خیلی برایم جالب خواهد بود که به مطالعه افراد بپردازم.» این کار داشت نتیجه می‌داد. او
مثل این بود که باید به خودش در آینه نگاه کند تا مطمئن شود که
داشت باز می‌شد. او مرتب صحبت می‌کرد و می‌کرد: «و می‌توانی با بعضی
آنچه که تازه اتفاق افتاده واقعی بود.
کسی برخورد کنی یا با کسی وارد رابطه شوی و بگویی، 'آیا آنها الان به من دروغ
«وای،» نفس‌نفس‌زنان گفت. «من این کار را کردم.» او مثل دختر بچه‌ای بود که برای اولین بار بریتنی
می‌گویند؟ ای‌خدای من.»
اسپیرز را می‌بیند. او طرفدار خودش بود.
وقت آن رسیده بود که اسلحه‌های سنگین را به میدان بیاوریم.
«فقط می‌دانستم که عدد هفت است!» او اعلام کرد و به سوی
«چیزی واقعاً جالب نشان‌ات می‌دهم و بعد به مصاحبه برمی‌گردیم,» گفتم، با افزودن محدودیت زمانی برای تاثیرگذاری بهتر. «یک آزمایش خواهد بود. می‌خواهم سعی کنم چیزی را که در افکارت است حدس بزنم.»
مبل بازگشت. البته که او می‌دانست. آن اولین تردستی بود که از میستری یاد گرفته بودم: اگر کسی بین یک تا ده یک عدد را به صورت تصادفی انتخاب کند، هفتاد درصد مواقع، مخصوصاً اگر تصمیم‌گیری‌اش را عجله کنید، آن عدد هفت خواهد بود.
بنابراین، بله، او را فریب داده بودم. اما اعتماد به نفس او به یک تقویت خوب نیاز داشت.
«می‌بینی،» به او گفتم. «تو از قبل تمام پاسخ‌ها را در درون خودت داری. فقط جامعه تو را مجبور به زیاد فکر کردن می‌کند.» واقعاً به این معتقد بودم.
«مصاحبه باحالی بود!» او اظهار داشت. «این مصاحبه را دوست دارم! این بهترین مصاحبه زندگی من بوده است!»
سپس او به سمت من چهره‌اش را چرخاند، به چشمان من نگاه کرد و پرسید، «می‌توانیم ضبط‌صوت را متوقف کنیم؟»
برای پانزده دقیقه بعد، درباره معنویت و نوشتن و زندگی‌های‌مان صحبت کردیم. او فقط یک دختر کوچولوی گمشده بود که در حال گذراندن بلوغ احساسی دیرهنگامی بود. او به دنبال چیزی واقعی برای نگه داشتن بود، چیزی عمیق‌تر از شهرت پاپ و چاپلوسی دستیارانش. من ارزش‌ها را نشان داده بودم و اکنون به مرحله ایجاد رابطه می‌رفتیم. شاید میستری درست می‌گفت: همه روابط انسانی تابع یک فرمول هستند.
رابطه برابر است با اعتماد به علاوه راحتی.
اما، من وظیفه‌ای برای انجام داشتم. ضبط‌صوت را روشن کردم و سوالاتی را که در ابتدای مصاحبه به او داده بودم و همه سوالات دیگر را هم پرسیدم. این بار او پاسخ‌های واقعی به من داد، پاسخ‌هایی که می‌توانستم چاپ کنم.
وقتی ساعت تمام شد، ضبط‌صوت را متوقف کردم.
«می‌دانی،» بریتنی گفت. «همه چیز به دلیلی اتفاق می‌افتد.»
«من واقعاً به این باور دارم،» به او گفتم.
«من هم همین‌طور.» او شانه‌ام را لمس کرد و لبخندی بزرگ در جریان چهره‌اش پخش شد. «دوست دارم شماره‌ها را با هم عوض کنیم.»
به آرامی برگه را برگرداند، طوری که انگار می‌ترسید نگاه کند، سپس آن را به سطح چشم آورد و یک عدد هفت بزرگ را دید که مستقیماً به او نگاه می‌کرد.
او جیغ کشید، از روی مبل پرید و به سوی آینه هتل دوید. دهانش باز ماند وقتی که در انعکاس به چشم‌هایش نگاه کرد.
«ای خدای من،» به انعکاسش گفت. «من این کار را کردم.»